\begin{table*}[t]\centering\footnotesize
\caption{Research questions, systematization findings and where each is established. Experiments test the comparison rules; the SoK contribution is the record-level synthesis.}\label{tab:rq-findings}
\begin{tabularx}{\textwidth}{@{}>{\raggedright\arraybackslash}p{.15\textwidth}Y>{\raggedright\arraybackslash}p{.29\textwidth}@{}}\toprule
Question & Finding & Evidence\\\midrule
RQ1: targets and observables & A common signal can serve confidentiality, attribution, evasion or authorized diagnostics; polarity changes the protected interest. Signal generation clusters at serving (32 families) and orchestration (19), whereas observation clusters at application, hardware and transport. & Complete 58-family core map (\S\ref{sec:corpus-map}, Table~\ref{tab:core-family-map}).\\
RQ2: origin versus access & Packet-size cover does not bind call order; plaintext re-emission creates a fresh size trace; shared call projections do not establish unchanged state. Within one paper, 7 of 15 same-family record pairs measure different primitives, targets or cost objects. & 39-record ledger (\S\ref{sec:defenses}); NetEcho/SplitDecoy and WRA/AHB lenses (\S\ref{sec:discussion}); Exps~1--2.\\
RQ3: portability & Explicit holdouts exist for task/input in 13 of 82 families, for model in three, and for neither hardware nor time in any; input separation does not certify a changed generator. & Six-axis matrix (\S\ref{sec:evaluation}, Appendix~\ref{app:holdouts}); calibration-transfer matrix in Exp~3.\\
RQ4: protection, state and cost & Inference is bound far more often than reuse: matched adaptation 2/11 applicable families, utility at the protected operating point 3/17, cost 11/19. & Nine-item, 20-family audit (\S\ref{sec:defenses}); Exps~3--4 leakage/state/cost contrasts.\\\bottomrule
\end{tabularx}\end{table*}
